% from ega0-19-fr.tex
\oldpage[0\textsubscript{IV-1}]{90}
Soient $C$ un anneau topologique (commutatif), $V$ un $C$-module topologique,
et considérons l'algèbre symétrique
$\mathbf S_C^*(V)=\bigoplus_{n\geq0}\mathbf S_C^n(V)$, que nous allons munir
canoniquement d'une \emph{topologie linéaire}, compatible avec sa structure de
$C$-algèbre. Pour cela, considérons un sous-module ouvert $U$ de $V$, et
l'\emph{idéal gradué} $U.\mathbf S_C^*(V)$ qu'il engendre dans
$\mathbf S_C^*(V)$~; nous prendrons comme système fondamental de voisinages de
$0$ dans $\mathbf S_C^*(V)$ l'ensemble des sommes
$\mathfrak a.\mathbf S_C^*(V)+U.\mathbf S_C^*(V)$, où $\mathfrak a$ (resp.
$U$) parcourt un système fondamental d'idéaux ouverts (resp. de sous-modules
ouverts) de $C$ (resp. $V$). On notera que si la topologie de $V$ est moins
fine que la topologie déduite de celle de $C$, on peut se borner aux couples
$(\mathfrak a,U)$ tels que $\mathfrak aV\subset U$, donc
$\mathfrak a.\mathbf S_C^*(V)+U.\mathbf S_C^*(V)
=\mathfrak a.1_C+U.\mathbf S_C^*(V)$~; dans ce cas la topologie induite sur
chacun des $\mathbf S_C^n(V)$ pour $n\geq1$ a pour système fondamental de
voisinages de $0$ les $U.\mathbf S_C^{n-1}(V)$, où $U$ parcourt les
sous-modules ouverts de $V$~; en particulier, sur $V=\mathbf S_C^1(V)$, elle
coïncide avec la topologie donnée (en général elle est moins fine que cette
dernière). En outre, dans tous les cas, l'\emph{algèbre topologique}
$\mathbf S_C^*(V)$ ainsi définie est, pour les catégories des $C$-modules
topologiques et des $C$-algèbres topologiques, la solution du même problème
universel que pour les catégories de $C$-modules et de $C$-algèbres. En effet,
soit $E$ une $C$-algèbre topologique, $u:V\to E$ un homomorphisme de
$C$-modules, $\mathbf S(u)=\widetilde u$ son extension canonique à
$\mathbf S_C^*(V)$. Supposons $u$ continu~; alors, si $\mathfrak b$ est un
idéal ouvert de $E$, son image réciproque $U=u^{-1}(\mathfrak b)$ est un
sous-$C$-module ouvert de $V$, et l'image par $\widetilde u$ de
$U.\mathbf S_C^*(V)$ est donc contenue dans $\mathfrak b$~; comme par ailleurs
il existe un idéal ouvert $\mathfrak a$ de $C$ tel que
$\mathfrak a.E\subset\mathfrak b$, donc
$\widetilde u(\mathfrak a.\mathbf S_C^*(V))\subset\mathfrak b$, cela prouve
que $\widetilde u$ est continu. Réciproquement, si $\widetilde u$ est continu,
et si $\mathfrak b$ est un idéal ouvert de $E$, il existe un sous-module ouvert
$U$ de $V$ tel que $\widetilde u(U.\mathbf S_C^*(V))\subset\mathfrak b$, et en
particulier $\widetilde u(U.1_C)\subset\mathfrak b$, c'est-à-dire
$u(U)\subset\mathfrak b$, donc $u$ est continu. Rappelons en outre que l'on a
un isomorphisme canonique de $C$-modules topologiques (discrets)
\begin{equation}
\label{0.19.5.1.1-fr}
  \mathbf S_C^n(V)/U.\mathbf S_C^{n-1}(V)
  \xrightarrow{\sim}\mathbf S_C^n(V/U).
  \tag{19.5.1.1}
\end{equation}
\end{env}

\begin{env}[19.5.2]
\label{0.19.5.2-fr}
Soient $A$ un anneau topologique, $B$ une $A$-algèbre topologique
(commutative), $\mathfrak J$ un idéal de $B$ (non nécessairement ouvert ni
fermé)~; dans tout ce qui suit, on munit $\mathfrak J$ et les
$\mathfrak J^n$ ($n\geq2$) de la topologie \emph{induite} par celle de $B$,
les quotients $C=B/\mathfrak J$,
$\mathfrak J^n/\mathfrak J^{n+1}=\gr_{\mathfrak J}^n(B)$ de la topologie
\emph{quotient}, de sorte que les $\mathfrak J^n/\mathfrak J^{n+1}$ sont des
$C$-modules topologiques~; l'injection canonique
$\mathfrak J/\mathfrak J^2\to\gr_{\mathfrak J}^*(B)$ se prolonge en un
homomorphisme (d'abord non topologique) de $C$-algèbres
\[
  \varphi:\mathbf S_C^*(\mathfrak J/\mathfrak J^2)
  \longrightarrow\gr_{\mathfrak J}^*(B)
\]
qui pour chaque $n$ donne un homomorphisme \emph{surjectif} de $C$-modules
\begin{equation}
\label{0.19.5.2.1-fr}
  \varphi_n:\mathbf S_C^n(\mathfrak J/\mathfrak J^2)
  \longrightarrow\mathfrak J^n/\mathfrak J^{n+1}
  =\gr_{\mathfrak J}^n(B).
  \tag{19.5.2.1}
\end{equation}

